\chapter{Introduction}

\section{First Section}
First paragraph \cite{Doe}.

Second paragraph.  Sample reference to \Fig{fig:above-below}.  Just some extra text to see how the text
wrapping looks in a compiled pdf document.

% Single Figure
\begin{figure}
  \centering
  \includegraphics[trim=0pt 0pt 0pt 0pt, clip=true,
  keepaspectratio=true, width=0.5\textwidth]{universe.jpg} \\
  \caption{Sample of a single image.}\label{fig:single2}
\end{figure}

Just some extra text to see how the text wrapping looks in a
compiled pdf document. Just some extra text to see how the text
wrapping looks in a compiled pdf document. Just some extra text to
see how the text wrapping looks in a compiled pdf document. Just
some extra text to see how the text wrapping looks in a compiled
pdf document. Just some extra text to see how the text wrapping
looks in a compiled pdf document.
\newpage
\begin{enumerate}
	\item item1 item1 item1 item1 item1 item1 item1 item1 item1 item1 item1 item1 item1 item1 item1 item1 item1 item1 
	\item item2 item2 item2 item2 item2 item2 item2 item2 item2 item2 item2 item2 item2 item2 item2 item2 item2 item2 
	\item item3 item3 item3 item3 item3 item3 item3 item3 item3 item3 item3 item3 item3 item3 item3 item3 item3 item3 
\end{enumerate}

\begin{enumerate}
	\packed
	\item item1 item1 item1 item1 item1 item1 item1 item1 item1 item1 item1 item1 item1 item1 item1 item1 item1 item1 
	\item item2 item2 item2 item2 item2 item2 item2 item2 item2 item2 item2 item2 item2 item2 item2 item2 item2 item2 
	\item item3 item3 item3 item3 item3 item3 item3 item3 item3 item3 item3 item3 item3 item3 item3 item3 item3 item3 
\end{enumerate}

\begin{enumerate}
	\packed
	\item One
	\item Two
	\item Three
	\begin{enumerate}
		\item test
		\item best
	\end{enumerate}
\end{enumerate}

\begin{itemize}
	\packed
	\item A
	\item B
	\item C
\end{itemize}

Just some extra text to see how the text wrapping looks in a
compiled pdf document. Just some extra text to see how the text
wrapping looks in a compiled pdf document. Just some extra text to
see how the text wrapping looks in a compiled pdf document. Just
some extra text to see how the text wrapping looks in a compiled
pdf document.

% Side-by-side subfigure
\begin{figure}
     \centering
      \subfloat[Case I]{\includegraphics[width=0.1\textwidth]{CUlogo.eps}}\hfil
      \subfloat[Case II]{\includegraphics[width=0.4\textwidth]{DOE.png}}
    \caption{Samble side-by-side subfigures} \label{fig:side-by-side}
\end{figure}

Just some extra text to see how the text wrapping looks in a
compiled pdf document. Just some extra text to see how the text
wrapping looks in a compiled pdf document.

% Above-below subfigure
\begin{figure}
    \centering
    \subfloat[Case I]{\includegraphics[width=0.1\textwidth]{CUlogo.eps}} \hspace{24pt}
    \subfloat[Case II]{\includegraphics[width=0.1\textwidth]{CUlogo.png}}
    \caption{Samble above-below subfigures} \label{fig:above-below}
\end{figure}

\section{Section Demonstrating Tables}
First paragraph.  This is a sample reference to \Tab{tab:sample}.

\begin{table}
  \centering
  \caption{Sample table. With an extra long caption to test how captions will wrap.}\label{tab:sample}
    \vspace{6pt} % Required to get proper spacing between caption and table
  \begin{tabular}{c c c}
    \hline
    Label 1 & Label 2 & Label 3 \\
    \hline\hline
    value 1 & $x_1$ & $y_1$ \\
    value 2 & $x_2$ & $y_2$ \\
    value 3 & $x_3$ & $y_3$ \\
    \hline
  \end{tabular}
\end{table}

New paragraph.

\begin{equation}\label{eqn:newton_law}
    F = Ma
\end{equation}

Paragraph referencing an \eqref{eqn:newton_law}.
